
\subsubsection*{5.1.4}
\phantomsection
\label{sga3:VIIB:5.1.4}

Let \((A,J)\) be an object of \(\mathcal C\), and let
\[
  (A_n)=\operatorname{Gr}_J(A)
\]
be the associated object of \(\mathbf{Alpg}/k\); thus
\[
  A_n=\overline{J^n}/\overline{J^{n+1}}
  \qquad(n\geq0).
\]
It is clear that the algebra
\[
  \mathscr A=\prod_{n\geq0}A_n
\]
is generated by \(A_0\) and \(A_1\). In other words, for
\(n\geq1\), the multiplication map
\[
  A_1\widehat\otimes_{A_0}A_1
  \widehat\otimes_{A_0}\cdots
  \widehat\otimes_{A_0}A_1
  \longrightarrow A_n
  \tag{1}
  \label{sga3:VIIB:5.1.4:eq1}
\]
is surjective.

Suppose further that \((A,J)\) is a cogroup of \(\mathcal C\)
and that \(A\) and the quotients
\[
  \overline{J^n}/\overline{J^{n+1}}
\]
are topologically flat over \(k\). By \SGARefLink{sga3:VIIB:5.1.2:corollary}{Corollary~5.1.2},
\(\operatorname{Gr}_J(A)\) is then a cogroup of
\(\mathbf{Alpg}/k\). Hence, by \SGARefLink{sga3:VIIB:5.1.3}{5.1.3}, if one puts
\[
\begin{split}
  A'_n
  =
  \biggl\{
    x\in\overline{J^n}/\overline{J^{n+1}}
    \ \biggm|\
    \Delta(x)-x\widehat\otimes1
    \in
    \bigoplus_{i=1}^{n}
    &\bigl(
      \overline{J^{\,n-i}}/\overline{J^{\,n-i+1}}
    \bigr)
    \widehat\otimes_k
    \\
    &\bigl(
      \overline{J^i}/\overline{J^{i+1}}
    \bigr)
  \biggr\},
\end{split}
\tag{2}
\label{sga3:VIIB:5.1.4:eq2}
\]
then the map
\[
  A'_n\widehat\otimes_kA_0\longrightarrow A_n,
  \qquad
  a'_n\widehat\otimes a_0\longmapsto a'_na_0,
\]
is an isomorphism in \(\mathbf{Alpg}/k\).\footnote{%
\label{sga3:VIIB:nde:149}Let \(G\), \(H\), and
\(\widehat G_H\) be the formal \(k\)-groups corresponding to
\(A\), \(A_0=A/J\), and \(\operatorname{Gr}_J(A)\),
respectively. Then
\(\tau:A_0\hookrightarrow\operatorname{Gr}_J(A)\) corresponds
to a retraction
\(\pi:\widehat G_H\to H\) of the inclusion
\(H\hookrightarrow\widehat G_H\), and the preceding assertion
means that \(\widehat G_H\) is the semidirect product of
\(N=\ker(\pi)\) by \(H\).}

\phantomsection
\label{sga3:VIIB:5.1.4:eq2-source-tensor-note}
\noindent\textit{Source note.}\enspace
In the displayed element-wise description of this map, the
French PDF prints
\(a'_n\widehat\otimes a'_n\mapsto a'_na_0\).
The domain and the product require
\(a'_n\widehat\otimes a_0\mapsto a'_na_0\), which is used here.

Let \(I_0\) denote the augmentation ideal of \(A_0\). It follows
from \textup{\SGARefLink{sga3:VIIB:5.1.4:eq1}{(1)}} and the following commutative diagram, in which
\[
  A_1'{}^{\widehat\otimes n}
  =
  \underbrace{
    A'_1\widehat\otimes_k\cdots\widehat\otimes_kA'_1
  }_{n\text{ factors}},
\]
that:
\begin{center}
  \phantomsection
  \label{sga3:VIIB:5.1.4:degree-one-generation-diagram}
  \begin{tikzcd}[row sep=3em,column sep=2.1em]
A_1\widehat\otimes_{A_0}\cdots\widehat\otimes_{A_0}A_1
  \arrow[d,"m"']
  & A_1'{}^{\widehat\otimes n}\widehat\otimes_k A_0
      \arrow[l,"{\sim}"']
      \arrow[d,"{m'\widehat\otimes\operatorname{id}}"']
  & \bigl(A_1'{}^{\widehat\otimes n}\widehat\otimes_k I_0\bigr)
      \oplus A_1'{}^{\widehat\otimes n}
      \arrow[l,"{\sim}"']
      \arrow[d,"{(m'\widehat\otimes\operatorname{id})\oplus m'}"] \\
A_n
  & A'_n\widehat\otimes_k A_0 \arrow[l,"{\sim}"']
  & (A'_n\widehat\otimes_k I_0)\oplus A'_n
      \arrow[l,"{\sim}"']
\end{tikzcd}

\end{center}
the multiplication map
\[
  m':
  A'_1\widehat\otimes_k\cdots\widehat\otimes_kA'_1
  \longrightarrow A'_n
  \tag{3}
  \label{sga3:VIIB:5.1.4:eq3}
\]
is surjective. In other words, the profinite \(k\)-algebra
\[
  \mathscr A'=\prod_{n\geq0}A'_n
\]
is generated by its degree-one terms.
